\subsection{R\'enyi 有效支撑、单圆 Lie 可见相位与 Dirichlet 门槛障碍}\label{subsec:conclusion-renyi-effective-support-circlephase-dirichlet-obstruction}
沿用定理 \ref{thm:fold-collision-spectrum}、推论 \ref{cor:fold-collision-scale-diverges}、推论 \ref{cor:real-input-40-ground-entropy}、定理 \ref{thm:conclusion-cdim-finite-rank-extension}、定理 \ref{thm:real-input-40-det-uv-2plus8}、推论 \ref{cor:real-input-40-zeta-uv-sqrtv-eigs}、定理 \ref{thm:xi-time-part68-entropy-deficit-triple-identification}、定理 \ref{thm:cdim-star-z1s-dual-extension}、定理 \ref{thm:xi-localized-connected-lie-quotient-circle-rigidity} 与命题 \ref{prop:real-input-40-output-dirichlet-nonrh} 的记号。前文已经分别在 bin-fold 碰撞谱、real-input \(40\) 的零温饱和层、二维加权 rigid/core 分裂、全阶 R\'enyi 缺损以及局部化 solenoid 的 Lie 可见层上给出闭式刻画。把这些接口并置之后，还可继续抽取出如下五条直接后果。它们共同表明：连续 Lie 可见部分始终压缩到单圆相位，而离散可恢复复杂度、二阶有效支撑与 Dirichlet 慢模障碍则保留不可消去的正速率。

\begin{theorem}[二阶 R\'enyi 有效支撑的零密度塌缩]\label{thm:conclusion-renyi2-effective-support-zero-density-collapse}
记
\[
M_m:=F_{m+2},
\qquad
p_m(r):=\frac{d_m(r)}{2^m},
\qquad
\mathsf{Col}_m:=\sum_{r\in \ZZ/(M_m\ZZ)}p_m(r)^2,
\]
并定义二阶 R\'enyi 熵与对应的有效支撑数
\[
H_2(p_m):=-\log\sum_{r}p_m(r)^2,
\qquad
N_{\mathrm{eff}}^{(2)}(m):=\exp\!\bigl(H_2(p_m)\bigr).
\]
则有精确恒等式
\[
N_{\mathrm{eff}}^{(2)}(m)=\frac{1}{\mathsf{Col}_m},
\]
并且
\[
\frac{N_{\mathrm{eff}}^{(2)}(m)}{M_m}
=
\frac{1}{M_m\mathsf{Col}_m}
\longrightarrow 0,
\qquad
\log M_m-H_2(p_m)=\log\!\bigl(M_m\mathsf{Col}_m\bigr)\longrightarrow +\infty.
\]
换言之，bin-fold 折叠分布在二阶碰撞几何下所占据的有效支撑，相对于整个 Fibonacci 窗口只形成零密度子集。
\end{theorem}
\leanverified{paper\_conclusion\_renyi2\_effective\_support\_zero\_density\_collapse}
\begin{proof}
由定理 \ref{thm:fold-collision-spectrum}，
\[
\sum_r\left(p_m(r)-\frac{1}{M_m}\right)^2
=
\mathsf{Col}_m-\frac{1}{M_m}.
\]
展开左端并利用 \(\sum_r p_m(r)=1\)，便得
\[
\sum_r p_m(r)^2=\mathsf{Col}_m.
\]
因此
\[
H_2(p_m)=-\log\mathsf{Col}_m,
\qquad
N_{\mathrm{eff}}^{(2)}(m)=e^{H_2(p_m)}=\frac{1}{\mathsf{Col}_m}.
\]
另一方面，推论 \ref{cor:fold-collision-scale-diverges} 给出
\[
M_m\mathsf{Col}_m\longrightarrow +\infty.
\]
代回即得
\[
\frac{N_{\mathrm{eff}}^{(2)}(m)}{M_m}
=
\frac{1}{M_m\mathsf{Col}_m}
\longrightarrow 0,
\qquad
\log M_m-H_2(p_m)=\log(M_m\mathsf{Col}_m)\longrightarrow +\infty.
\]
证毕。
\end{proof}

\begin{theorem}[一圆 Lie 可见度与正地基层熵的严格分离]\label{thm:conclusion-one-circle-visible-phase-positive-ground-entropy-separation}
沿用推论 \ref{cor:real-input-40-ground-entropy} 的记号。考虑任意精确解码机制：在已知路径属于最大输出密度层
\[
k=k_{\max}(n)
\]
的条件下，再辅以 \(B_n\) 比特侧信息，要求无误恢复该层中的每一条长度 \(n\) 可行路径。则必有
\[
B_n\ge \log_2 A_{n,k_{\max}(n)}
=
n\log_2 c+o(n),
\]
从而
\[
\liminf_{n\to\infty}\frac{B_n}{n}\ge \log_2 c>0.
\]
另一方面，对任意有限素数集 \(S\) 与局部化数系
\[
G_S:=\ZZ[S^{-1}],
\]
都有
\[
\cdim_{\mathrm{fr}}(G_S)=1.
\]
因此，在有限秩圆维口径下，连续 Lie 可见层始终只保留一条圆相位，而极端饱和层的精确恢复仍保留严格正的离散熵率 \(\log_2 c\)。
\end{theorem}
\begin{proof}
最大输出密度层中的对象总数恰为 \(A_{n,k_{\max}(n)}\)。若只允许 \(B_n\) 比特侧信息而又要求对该层中的每一条路径都无误恢复，则不同路径必须对应不同证书，故证书数至多 \(2^{B_n}\) 时必有
\[
2^{B_n}\ge A_{n,k_{\max}(n)}.
\]
取对数即得
\[
B_n\ge \log_2 A_{n,k_{\max}(n)}.
\]
再由推论 \ref{cor:real-input-40-ground-entropy}，
\[
\lim_{n\to\infty}\frac{1}{n}\log A_{n,k_{\max}(n)}=\log c,
\]
于是
\[
\log_2 A_{n,k_{\max}(n)}=n\log_2 c+o(n),
\]
从而
\[
\liminf_{n\to\infty}\frac{B_n}{n}\ge \log_2 c>0.
\]

另一方面，\(G_S=\ZZ[S^{-1}]\) 为无挠有限秩阿贝尔群，且
\[
\operatorname{rank}_{\QQ}(G_S)=1.
\]
由定理 \ref{thm:conclusion-cdim-finite-rank-extension}，便得到
\[
\cdim_{\mathrm{fr}}(G_S)=1.
\]
证毕。
\end{proof}

\begin{proposition}[二维碰撞--输出谱中的刚性正交片]\label{prop:conclusion-rigid-orthogonal-slice-collision-output-spectrum}
沿用定理 \ref{thm:real-input-40-det-uv-2plus8} 的二维加权行列式
\[
\Delta(z;u,v)=(vz^2-1)P_8(z;u,v),
\]
并在 \(v>0\) 上取正的刚性谱支
\[
\lambda_{\mathrm{rig}}(u,v):=\sqrt v.
\]
定义相应谱压
\[
P_{\mathrm{rig}}(u,v):=\log\lambda_{\mathrm{rig}}(u,v)=\frac12\log v.
\]
则对任意整数 \(a\ge 1\)、\(b\ge 0\)，都有
\[
\partial_u^{\,a}\partial_v^{\,b}P_{\mathrm{rig}}(u,v)=0.
\]
因此，刚性谱支对“碰撞计数--输出计数”联合累积量的全部 \(u\)-敏感分量均无贡献；全部非平凡耦合只能由非刚性 core 因子 \(P_8(z;u,v)\) 承载。
\end{proposition}
\begin{proof}
由推论 \ref{cor:real-input-40-zeta-uv-sqrtv-eigs}，二维加权模型的两条 rigid branches 为
\[
\lambda_\pm=\pm\sqrt v.
\]
取正支即得
\[
P_{\mathrm{rig}}(u,v)=\log\sqrt v=\frac12\log v,
\]
其右端与 \(u\) 无关。故凡至少含有一次 \(u\)-微分的混合导数必为零，即
\[
\partial_u^{\,a}\partial_v^{\,b}P_{\mathrm{rig}}(u,v)=0
\qquad (a\ge 1,\ b\ge 0).
\]
证毕。
\end{proof}

\begin{corollary}[R\'enyi--普适缺损常数与 oracle 速率下界的同一化]\label{cor:conclusion-renyi-universal-defect-oracle-rate-identification}
沿用定理 \ref{thm:xi-time-part68-entropy-deficit-triple-identification} 的记号。对任意有限 \(q>0\)（其中 \(q=1\) 按 Shannon 熵解释），定义
\[
\delta_q(m):=m\log 2-H_q(\mu_m).
\]
则有与 \(q\) 无关的极限
\[
\lim_{m\to\infty}\frac{\delta_q(m)}{m}
=
\log\frac{2}{\varphi}.
\]
若一列比特预算 \(B_m\) 满足
\[
\mathrm{Succ}_m^{\mathrm{bin}}(B_m)\longrightarrow 1,
\]
则必有
\[
\liminf_{m\to\infty}\frac{B_m}{m}
\ge
\log_2\frac{2}{\varphi}.
\]
因此，同一个常数 \(\log(2/\varphi)\) 同时给出全部有限阶 R\'enyi 缺损的广延速率与无结构精确反演的普适 oracle 比特率下界。
\end{corollary}
\begin{proof}
定理 \ref{thm:xi-time-part68-entropy-deficit-triple-identification} 已给出：对任意 \(q>0\)，
\[
\lim_{m\to\infty}\frac{1}{m}\bigl(m\log 2-H_q(\mu_m)\bigr)=\log\frac{2}{\varphi}.
\]
这正是第一条。

若 \(\mathrm{Succ}_m^{\mathrm{bin}}(B_m)\to 1\)，则对任意固定 \(\varepsilon\in(0,1)\)，当 \(m\) 充分大时都有
\[
\mathrm{Succ}_m^{\mathrm{bin}}(B_m)\ge 1-\varepsilon.
\]
再次应用定理 \ref{thm:xi-time-part68-entropy-deficit-triple-identification}，即可得到
\[
\liminf_{m\to\infty}\frac{B_m}{m}\ge \log_2\frac{2}{\varphi}.
\]
证毕。
\end{proof}

\begin{theorem}[单圆 Lie 可见因子上的平方根门槛障碍]\label{thm:conclusion-circle-lie-factor-square-root-threshold-obstruction}
\leanverified{paper\_conclusion\_circle\_lie\_factor\_square\_root\_threshold\_obstruction}
设 \(S\subset\mathbb P\) 为有限素数集，并记
\[
G_S:=\ZZ[S^{-1}],
\qquad
\Sigma_S:=\widehat{G_S}.
\]
则 \(\Sigma_S\) 落在短正合列
\[
0\longrightarrow \prod_{p\in S}\ZZ_p\longrightarrow \Sigma_S\xrightarrow{\pi_S}\TT\longrightarrow 0
\]
之中，并且 \(\Sigma_S\) 的任意非平凡连通紧 Lie 商都同构于 \(\TT\)。

另一方面，对输出 \(1\) 计数模 \(m\) 的 Dirichlet 扭曲，定义
\[
\theta_m:=\max_{1\le j\le m-1}\frac{\log\rho_{m,j}}{\log\lambda},
\qquad
\lambda=\varphi^2.
\]
则存在某个模数 \(m\ge 2\) 使
\[
\theta_m>\frac12.
\]
因此，RH 型的平方根门槛不可能在全部圆相位上统一成立。更精确地，若某个扭曲相位 \(u=\omega_m^j\in\TT\) 满足
\[
\rho_{m,j}=\lambda^{\theta_{m,j}},
\qquad
\theta_{m,j}>\frac12,
\]
则对任意
\[
0<\varepsilon<\theta_{m,j}-\frac12
\]
都不可能在该相位上成立界
\[
O\!\bigl(\lambda^{n(1/2+\varepsilon)}\bigr).
\]
换言之，平方根障碍首先出现在这一唯一类型的连通 Lie 可见商 \(\TT\) 上，而不是先由 profinite 核 \(\prod_{p\in S}\ZZ_p\) 触发。
\end{theorem}
\begin{proof}
短正合列正是定理 \ref{thm:cdim-star-z1s-dual-extension} 的内容；而由定理 \ref{thm:xi-localized-connected-lie-quotient-circle-rigidity}，\(\Sigma_S\) 的任意非平凡连通紧 Lie 商都只能同构于 \(\TT\)。

另一方面，命题 \ref{prop:real-input-40-output-dirichlet-nonrh} 已证明：存在明确模数 \(m\) 使
\[
\theta_m>\frac12.
\]
故可取某个 \(j\) 使 \(\theta_{m,j}>\frac12\)，也即
\[
\rho_{m,j}=\lambda^{\theta_{m,j}}>\lambda^{1/2}.
\]
对任意 \(0<\varepsilon<\theta_{m,j}-1/2\)，都有
\[
\lambda^{\theta_{m,j}}>\lambda^{1/2+\varepsilon}.
\]
而同一命题已经指出，对应该扭曲通道的 Dirichlet/Fourier 误差不具备 RH 形的平方根指数，因而不可能在该相位上满足
\[
O\!\bigl(\lambda^{n(1/2+\varepsilon)}\bigr).
\]
最后，见证相位 \(u=\omega_m^j\) 本身就是圆周 \(\TT\) 上的 torsion 点，故障碍已在一圆 Lie 可见因子上显现。与此同时，当 \(S\neq\varnothing\) 时，扩张的严格不分裂性由定理 \ref{thm:xi-localized-basic-extension-strictly-nonsplit} 给出，这说明 \(\TT\) 只记录 Lie 可见通道，而不等同于整个 solenoid。证毕。
\end{proof}

\noindent
上述五条结果把 bin-fold 碰撞几何、极端饱和层的恢复复杂度、二维 rigid/core 分裂、全阶 R\'enyi 缺损以及局部化 solenoid 的 Lie 可见结构压缩为同一条分裂链：一方面，连续可见口径始终只保留一圆维；另一方面，二阶有效支撑的零密度、\(\log_2 c\) 的正恢复率以及 Dirichlet 慢模越过平方根门槛，又说明离散复杂度既不被单圆相位吞并，也不被纯 \(p\)-进核单独解释，而是在 circle quotient 与 arithmetic ledger 之间保持严格分工。

\endinput
